\chapter{Reference: testing programs}
\label{ch:refman-test}

\begin{aside}
The whole point of the Model–Update–View shape is that
testing is trivial. \code{update} is a pure function from
(Model, Msg) to (Model, Cmd). You don't need a terminal,
a fixture, or a mock; you need a unit test.
\end{aside}

\section{Working example}

\begin{lstlisting}[language=C++]
// counter.hpp
#pragma once
#include <maya/maya.hpp>

struct Counter {
    struct Model { int count = 0; };
    struct Inc {};
    struct Dec {};
    struct Quit {};
    using Msg = std::variant<Inc, Dec, Quit>;

    static Model init() { return {}; }

    static auto update(Model m, Msg msg)
        -> std::pair<Model, maya::Cmd<Msg>>
    {
        return std::visit(maya::overload{
            [&](Inc) { m.count++; return std::pair{m, maya::Cmd<Msg>{}}; },
            [&](Dec) { m.count--; return std::pair{m, maya::Cmd<Msg>{}}; },
            [](Quit) {
                return std::pair{Model{}, maya::Cmd<Msg>::quit()};
            },
        }, msg);
    }

    static maya::Element view(const Model& m) {
        return maya::dsl::text(m.count);
    }

    static auto subscribe(const Model&) -> maya::Sub<Msg> {
        return maya::dsl::key_map<Msg>({});
    }
};
\end{lstlisting}

\begin{lstlisting}[language=C++]
// counter_test.cpp
#include "counter.hpp"
#include <cassert>

int main() {
    auto m0 = Counter::init();
    assert(m0.count == 0);

    auto [m1, _1] = Counter::update(m0, Counter::Inc{});
    assert(m1.count == 1);

    auto [m2, _2] = Counter::update(m1, Counter::Inc{});
    assert(m2.count == 2);

    auto [m3, _3] = Counter::update(m2, Counter::Dec{});
    assert(m3.count == 1);

    return 0;
}
\end{lstlisting}

\section{What to test}

\begin{itemize}[leftmargin=*]
  \item \code{update} for every \code{Msg} variant.
  \item Invariants in the Model (counter never negative,
    selection always in range).
  \item Returned commands when relevant — assert
    \code{cmd.is\_quit()} on the Quit branch, for example.
\end{itemize}

\section{What \emph{not} to test}

\begin{itemize}[leftmargin=*]
  \item \code{view} pixel-by-pixel. Snapshot tests grow
    brittle and tell you nothing useful.
  \item The runtime itself. \maya{} has its own tests; trust
    them.
\end{itemize}

\section{Recap}

\begin{recap}
\begin{itemize}[leftmargin=*]
  \item Pure functions invite unit tests.
  \item Test \code{update}; assert on Model after each Msg.
  \item Skip the view.
\end{itemize}
\end{recap}

\section*{Workshop}

\begin{enumerate}[leftmargin=*]
  \item Write a property test: \code{Inc} then \code{Dec}
    is a no-op.
  \item Test that \code{Quit} always returns the quit
    command.
  \item Add a test for an edge case in your real app.
\end{enumerate}
